% ============================================================
%  CHAPTER 7 — WEB DASHBOARD
% ============================================================
\chapter{Web Analytics Dashboard}

\section{Technology Stack}

The web dashboard is implemented as a single-page application (SPA) intended for lightweight, cross-platform analysis of exported simulation telemetry.

\begin{table}[H]
\centering
\caption{Web Dashboard Technology Stack}
\label{tab:webstack}
\begin{tabularx}{\textwidth}{@{}lX@{}}
\toprule
\textbf{Technology} & \textbf{Role} \\
\midrule
React & Component-based UI for composing interactive views \\
Vite & Fast development server and production bundling \\
Tailwind CSS & Utility-first styling for consistent layout and spacing \\
Recharts & Charting primitives for time-series and distribution plots \\
\bottomrule
\end{tabularx}
\end{table}

\section{Data Ingestion}

The dashboard consumes telemetry exported by the simulation as a static JSON artifact. Each record contains packet-level fields such as packet size, inter-arrival time, and observed flow duration. This approach keeps the web layer decoupled from the simulator while enabling repeatable analysis.

\section{UI Summary}

The UI is organized around a small number of operator-focused views:
\begin{itemize}
    \item \textbf{KPI tiles} summarizing packet counts and aggregate latency statistics.
    \item \textbf{Time-series charts} comparing observed latency against model estimates and short-horizon forecasts.
    \item \textbf{Distribution plots} for identifying skew and outliers.
    \item \textbf{Topology view} representing the sensor--gateway--cloud path for contextual interpretation.
\end{itemize}

The design goal is clarity: consistent scales, readable labels, and low visual noise so that abnormal behavior is easy to spot.
